\label{chap.6}
\section{Method}
Every system sees the same events from the 2,304 held-out sessions. A session counts as detected if at least one of its events is held. Recall is the share of attack sessions that were detected, FPR is the share of harmless sessions that were wrongly held, and accuracy covers both. Confidence intervals are 95\% Wilson intervals~\cite{wilson1927}. MAPIS and the stateless model were trained five times with seeds 42, 2, 3, 4 and 5, and we report mean $\pm$ standard deviation. The other systems do not depend on a seed.

\section{Headline Results}
\begin{table}[H]
\centering\small
\begin{tabular}{lrrr}
\toprule
System & Recall (\%) & FPR (\%) & Accuracy (\%)\\
\midrule
\textbf{MAPIS} & \textbf{98.0 $\pm$ 0.6} & \textbf{0.9 $\pm$ 0.8} & \textbf{98.5}\\
Stateless DeBERTa (same data, event only) & 94.1 $\pm$ 0.9 & 33.2 $\pm$ 3.5 & 82.0\\
Llama Guard 3-1B & 86.1 & 90.3 & 52.3\\
ProtectAI DeBERTa-v3 injection v2 & 53.3 & 49.5 & 52.1\\
Prompt Guard 2 (86M) & 7.5 & 0.1 & 48.4\\
\bottomrule
\end{tabular}
\caption{Session-level results on 2,304 held-out sessions.}
\label{tab:head}
\end{table}

\begin{figure}[H]
\centering
\begin{tikzpicture}
\begin{axis}[ybar, bar width=9pt, width=13.5cm, height=6.4cm, ymin=0, ymax=105, ylabel={Percent}, symbolic x coords={MAPIS,Stateless,Llama Guard,ProtectAI,Prompt Guard 2}, xtick=data, x tick label style={font=\scriptsize}, legend style={at={(0.5,-0.18)},anchor=north,legend columns=2,font=\small}, nodes near coords, nodes near coords style={font=\tiny}, enlarge x limits=0.12]
\addplot[fill=maroon] coordinates {(MAPIS,98.0) (Stateless,94.1) (Llama Guard,86.1) (ProtectAI,53.3) (Prompt Guard 2,7.5)};
\addplot[fill=gray] coordinates {(MAPIS,0.9) (Stateless,33.2) (Llama Guard,90.3) (ProtectAI,49.5) (Prompt Guard 2,0.1)};
\legend{Recall, False-positive rate}
\end{axis}
\end{tikzpicture}
\caption{Recall and false-positive rate on 2,304 held-out sessions.}
\end{figure}

Llama Guard flags nearly every session and Prompt Guard 2 flags nearly none. ProtectAI flags about half of everything, and the stateless model catches most attacks only by also stopping a third of the harmless sessions. MAPIS keeps recall high and false alarms close to zero.

\section{Results by Set}
\begin{table}[H]
\centering\small
\begin{tabular}{lrr}
\toprule
Set & MAPIS recall (\%) & MAPIS FPR (\%)\\
\midrule
MAPIS-MultiHop & 100 & 3.0 $\pm$ 3.0\\
MAPIS-Adaptive (all six techniques) & 100 & 0.4\\
RealHarm: InjecAgent direct harm & 99.4 $\pm$ 0.7 & 0.0\\
RealHarm: ASB & 100 & 0.0\\
MAPIS-Decomposed & 100 & 0.0\\
MAPIS-Bench v1 & 100 & 0.0\\
BIPIA (never trained on) & 95.8 $\pm$ 1.4 & 1.4 $\pm$ 1.0\\
\bottomrule
\end{tabular}
\caption{MAPIS per benchmark set (mean over five seeds, default hold cut-off 0.50).}
\end{table}
For seed 42, recall by mechanism is 100\% for tool-output manipulation, cross-agent injection, goal hijacking and document poisoning, and 99.2\% for memory injection. By attack class, the weakest is instruction override at 92 to 96\%, and most of those misses are harmless-looking single-hop BIPIA injections.

\section{Matched False-Positive Budget}
At the default threshold the stateless model reaches 86.5\% recall on multi-hop attacks (five seeds), but with a false-positive rate of 51.1\%. In one seed it flags every authorized and designated twin. This is expected: an attack event and its authorized twin have exactly the same text, so a model that only sees the event has no way to tell them apart. For a fair comparison we limit the stateless model to at most 5\% false positives, pick the threshold on validation data and then apply it to the test data.
\begin{table}[H]
\centering\small
\begin{tabular}{lrr}
\toprule
Set & Stateless recall at $\le$5\% FPR (\%) & MAPIS recall (\%)\\
\midrule
MAPIS-MultiHop & 4.3 & 100\\
MAPIS-Adaptive & 6.1 & 100\\
InjecAgent direct harm & 88.5 & 99.4\\
ASB & 97.8 & 100\\
\bottomrule
\end{tabular}
\caption{Stateless recall at a matched false-positive budget (mean of five seeds).}
\end{table}
\begin{figure}[H]
\centering
\begin{tikzpicture}
\begin{axis}[ybar, bar width=14pt, width=11cm, height=6cm, ymin=0, ymax=110, ylabel={Recall (\%)}, symbolic x coords={Multi-hop,Adaptive}, xtick=data, legend style={at={(0.5,-0.15)},anchor=north,legend columns=2,font=\small}, nodes near coords, nodes near coords style={font=\scriptsize}, enlarge x limits=0.5]
\addplot[fill=maroon] coordinates {(Multi-hop,100) (Adaptive,100)};
\addplot[fill=gray] coordinates {(Multi-hop,4.3) (Adaptive,6.1)};
\legend{MAPIS, Stateless at $\le 5\%$ FPR}
\end{axis}
\end{tikzpicture}
\caption{Recall at a matched false-positive budget.}
\end{figure}
So at the same false-positive budget, MAPIS gains about 96 percentage points of recall over stateless detection on multi-hop attacks, while the target was 20. At the default threshold the plain gain is only 13.5 points (100 against 86.5), which is why we use the matched comparison. The thresholds picked on validation did not hold exactly on test (multi-hop FPR 7.7\% in seed 3, adaptive FPR 8.8\% in seed 2), so these stateless numbers are, if anything, generous to the baseline. On RealHarm the attack text is explicit, a message-level model does well even at zero false positives, and the gap is at most 11 points.

\section{The Evasive Attacker}
MAPIS-Adaptive assumes an attacker who knows how MAPIS works and tries to hide the planted item. Its test split uses held-out families and held-out evasion styles. Before the upgrade described in Section~\ref{sec:history}, the old model with the old matching caught 64.6\% of the adaptive attacks at 2.9\% FPR. With canonical matching and fragment matching the old model reached 100\%, but at 3.9\% FPR. After adversarial training the model catches every attack in all five seeds (576 sessions, 192 attacks), with 0.0\% FPR in four seeds and 1.8\% in one (mean 0.4\%). Table~\ref{tab:adaptive} breaks this down by technique.
\begin{table}[H]
\centering\small
\begin{tabular}{>{\raggedright\arraybackslash}p{2.9cm}>{\raggedright\arraybackslash}p{5.6cm}rr}
\toprule
Technique & Example & MAPIS & Stateless\\
\midrule
Obfuscated item & ``x at host dot example'', \texttt{hxxps://h[.]example} & 100 / 0 & 83.3 / 37.5\\
Split item & mailbox in one document, domain in another & 100 / 0 & 83.3 / 45.8\\
Hidden characters & zero-width characters, full-width @ & 100 / 0 & 87.5 / 43.8\\
Paraphrased claim & authority claim without the usual wording & 100 / 0 & 87.5 / 45.0\\
Long delay & eight hops between the plant and the action & 100 / 0 & 87.5 / 46.2\\
Combination & several techniques at once & 100 / 0 & 90.0 / 51.2\\
\bottomrule
\end{tabular}
\caption{MAPIS-Adaptive by technique, recall / FPR in \%, seed 42 (576 test sessions).}
\label{tab:adaptive}
\end{table}
The stateless model reaches 83 to 90\% recall here only by flagging 38 to 51\% of the twins. A long delay does not help the attacker against MAPIS, because a tainted item stays in the session record after its hop has left the four-hop history window.

\section{Public Guardrails by Set}
\begin{table}[H]
\centering\small
\begin{tabular}{lrrrrr}
\toprule
Set & MAPIS & Stateless & Llama Guard & ProtectAI & Prompt Guard 2\\
\midrule
All 2,304 sessions & 98.0 / 0.9 & 94.1 / 33.2 & 86.1 / 90.3 & 53.3 / 49.5 & 7.5 / 0.1\\
MAPIS-MultiHop & 100 / 3.0 & 86.5 / 51.1 & 100 / 100 & 70.8 / 63.9 & 0\\
MAPIS-Adaptive & 100 / 0.4 & 86.6 / 43.0 & 100 / 100 & 66.1 / 64.1 & 0\\
BIPIA (held-out part) & 95.8 / 1.4 & 94.5 / 2.2 & 72.5 / 65.5 & 44.7 / 56.5 & 1.0\\
\bottomrule
\end{tabular}
\caption{Recall / FPR in \% on the same sessions. MAPIS and the stateless model are means over five seeds; the public guardrails are pretrained and run once with their default thresholds. For Prompt Guard 2 the per-set cells give recall only.}
\end{table}
Llama Guard is a content-safety model and treats almost every agent session as unsafe. Prompt Guard 2 looks for explicit instructions such as ``ignore previous instructions'' and finds almost none in our attacks. ProtectAI flags attacks and harmless sessions at about the same rate. None of them can use the history of the session, so none can tell an attack from its twin.

\section{Ablation}
To see which part of MAPIS does the work, we remove one part at a time (seed 42, all 2,304 sessions).
\begin{table}[H]
\centering\small
\begin{tabular}{lrr}
\toprule
System & Recall (\%) & FPR (\%)\\
\midrule
\textbf{Full MAPIS} & \textbf{97.0} & \textbf{0.2}\\
Learned model only (no rule fusion) & 97.0 & 0.1\\
Without the provenance rule & 97.0 & 0.1\\
Same model without session memory & 94.7 & 33.6\\
Regex rules only & 12.5 & 0.0\\
\bottomrule
\end{tabular}
\caption{Ablation, seed 42.}
\end{table}
Removing the session memory raises the false-positive rate from 0.2\% to 33.6\%, and in the multi-hop twins the stateless model flags 89\% of the authorized and 77\% of the designated sessions, against 0.6\% and 0\% for MAPIS. Removing the fusion rules changes no decision: once the classifier reads the session signals in its input, it has learned what the rules add. The rules stay in the system as a safety net and because their reasons are what the operator sees in each alert.

\section{Calibrated Hold Cut-off}
With the default cut-off of 0.50, two of the five seeds go above 5\% false positives on multi-hop sessions (5.3 and 7.7\%). We therefore picked a cut-off on the validation split only, such that no dataset family goes above 2\% false positives there, and applied it to test. Across the five seeds, multi-hop FPR went from 1.2, 5.3, 0.0, 7.7 and 0.6\% to 0.6, 1.8, 0.0, 0.6 and 0.0\% (mean 3.0 to 0.6\%), recall stayed at 100\%, and no other set got worse. With this cut-off every seed meets the 5\% target.

\section{Latency and Cost}
On the RTX 3050 laptop GPU the classifier needs 63.0 ms per event when events come one at a time and 51.4 ms per event in batches. The DeBERTa-v3-base version, which we tried as a capacity comparison, needs 115 ms per event, still under 200 ms. The rule signals add only a few milliseconds.

\section{Capacity Comparison}
A DeBERTa-v3-base model (one seed, earlier data) reached 99.9\% recall and 2.0\% FPR, but 5.3\% FPR on multi-hop sessions. It gains some recall but costs false alarms and speed, so we kept the small model as the main one.

\section{Targets}
\begin{table}[H]
\centering\small
\begin{tabular}{p{4.2cm}p{6.6cm}p{3.2cm}}
\toprule
Target & Result & Met\\
\midrule
Accuracy $\ge$ 90\% & 98.5\% (5 seeds) & Yes\\
FPR $\le$ 5\% & $0.9\pm0.8$\% overall, multi-hop $3.0\pm3.0$\% & Yes (every seed $\le$1.8\% with the calibrated cut-off)\\
Latency $<$ 200 ms & 63 ms & Yes\\
Multi-hop recall gain $\ge$ 20 points & +96 points at matched FPR & Yes\\
\bottomrule
\end{tabular}
\caption{Project targets.}
\end{table}
